%============================================================
% TEMPLATE — ONE FULLY SOLVED GRAPH PROBLEM
% Copy to chapters/chNN_xxx/pK_xxx/probM_name.tex and fill in.
% Rules:
%   * complete C++14 in [style=cpp] (checked by tools/check_code.py),
%     fragments in [style=snippet], traces in [style=trace]
%   * box titles containing [ ] must be wrapped in braces:  [{title $d[v]$}]
%   * verify every number in traces and examples by running the code
%============================================================
\subsection{Problem N.K.M: Title}
\problemlabel{N.K.M}{ssec:chN-short-name}   % not \label: \ref must print the problem number

\problemcard{N.K.M}{Title}{Platform + id}{https://...}{\exmedium\ \cf{1600}}{one-line technique}

\subsubsection*{1. Problem statement}
% statement, constraints, sample table, a drawing of the sample graph (gv / ge / gde styles)

\subsubsection*{2. The graph model}
\begin{modelbox}[Graph Model --- Title]
\renewcommand{\arraystretch}{1.35}
\begin{tabularx}{\linewidth}{@{}>{\bfseries}l X@{}}
Vertices & what one vertex stands for (cell, state, number, person, \ldots) and how many there are \\
Edges & when $u \to v$ exists; directed or not; implicit or stored \\
Weights & none / 0-1 / non-negative / possibly negative / capacities \\
Question & reachability, shortest path, components, order, matching, \ldots{} in graph language \\
Size & $|V|$, $|E|$ in the worst case, and the time budget they allow \\
\end{tabularx}
\end{modelbox}
% picture: the problem on the left, the graph on the right

\subsubsection*{3. Key observations}
% the facts that make the algorithm work; why the obvious greedy / brute force fails (concrete counter-example)

\subsubsection*{4. Approaches}
\begin{bruteforcebox}
% the naive solution, its complexity, and the exact reason it is too slow (or wrong)
\end{bruteforcebox}
\begin{algorithmbox}
% the efficient algorithm: idea, complete code, complexitybox
\end{algorithmbox}
\begin{invariantbox}[Why it is correct]
% the invariant / lemma, with a short proof
\end{invariantbox}
\begin{alternativebox}
% a second correct algorithm (DFS vs BFS vs DSU, Kahn vs DFS order, ...) -- or say why none is natural
\end{alternativebox}
\begin{tricksbox}
% optional: iterative version, memory layout, constant factor, recursion depth
\end{tricksbox}

\subsubsection*{5. Trace}
\begin{tracebox}[{Sample 1}]
% the graph drawn after each important step (gvsrc, gvcur, gvseen, gvdone, gtree, gpath, gdist)
% plus a table of the data structure (queue, stack, dist[], parent[], DSU parents, ...)
\end{tracebox}

\subsubsection*{6. Test suite}
\begin{testbox}
% disconnected graph, single vertex, self-loops / multi-edges if allowed, largest input, counter-examples
\end{testbox}

\subsubsection*{7. Common mistakes}
% warningbox per mistake (wrong vs right)

\subsubsection*{8. Takeaways}
% 3-4 bullets, then the variants of this problem:
\begin{variantsbox}
\variant{Short name}{\exmedium\ \lc{Medium}\ \varlink{Judge title}{LeetCode 0}{https://leetcode.com/problems/slug/}}
  {what changes, one line}
  {model: what a vertex and an edge are now, at most two lines}
  {algorithm: what changes in the algorithm, at most two lines}
\variant{Already in the book}{\varin{https://leetcode.com/problems/slug-in-book/}}{...}{...}{...}
\variant{No judge}{\varnone}{...}{...}{...}
\end{variantsbox}
